\subsection{存在性与唯一性定理}

定义 2. --- 设 G 是局部紧群。G 上非零且左（分别地，右）不变的正测度称为 G 上的左（分别地，右）哈尔测度。

定理 1. --- 每个局部紧群上都存在左（分别地，右）哈尔测度，并且除相差一个常数因子外只有一个。

\begin{enumerate}
\def\labelenumi{\Alph{enumi})}
\tightlist
\item
  存在性。--- 置 \(\mathcal{K}(\mathrm{G}) = \mathcal{K}\)、\(\mathcal{K}_{+}(\mathrm{G}) = \mathcal{K}_{+}\)、\(\mathcal{K}_{+}^{*} = \mathcal{K}_{+} - \{0\}\)。若 C 是 G 的紧子集，则以 \(\mathcal{K}_{+}^{*}(\mathrm{C})\) 表示支集包含在 C 中的 \(f \in \mathcal{K}_{+}^{*}\) 的集合。对于 \(f \in \mathcal{K}\) 和 \(g \in \mathcal{K}_{+}^{*}\)，存在数 \(c_{1}, \ldots, c_{n} \geqslant 0\) 以及 G 中的元素 \(s_{1}, \ldots, s_{n}\)，使得 \(f \leqslant \sum_{i=1}^{n} c_{i} \gamma(s_{i}) g\)：这是因为，G 中存在非空开集 U，使得 \(\inf_{s \in \mathrm{U}} g(s) > 0\)，并且 \(f\) 的支集可以被 U 的有限个左平移覆盖。令 \((f : g)\) 为数 \(\sum_{i=1}^{n} c_{i}\) 的下确界，其中所有系统 \((c_{1}, \ldots, c_{n}, s_{1}, \ldots, s_{n})\) 由数 \(\geqslant 0\) 和 G 中的元素组成，并满足 \(f \leqslant \sum_{i=1}^{n} c_{i} \gamma(s_{i}) g\)。则：
\end{enumerate}

\begin{enumerate}
\def\labelenumi{(\roman{enumi})}
\item
  \(\left(\pmb{\gamma}(s)f:g\right) = (f:g)\) 对于 \(f\in \mathcal{K}\)、\(g\in \mathcal{K}_{+}^{*}\)、\(s\in G\) 成立；
\item
  \((\lambda f:g) = \lambda (f:g)\) 对于 \(f\in \mathcal{K}\) \(,g\in \mathcal{K}_{+}^{*},\lambda \geqslant 0\) 成立。
\item
  \(\left((f + f^{\prime}):g\right)\leqslant (f:g) + (f^{\prime}:g)\) 对于 \(f\in \mathcal{K},f^{\prime}\in \dot{\mathcal{K}},g\in \mathcal{K}_{+}^{*};\) 成立；
\item
  有 \((f:g)\geqslant (\sup f) / (\sup g)\)。
\end{enumerate}

\[
f \in \mathscr {K}, g \in \mathscr {K} _ {+} ^ {*}
\]

\begin{enumerate}
\def\labelenumi{(\alph{enumi})}
\setcounter{enumi}{21}
\tightlist
\item
  \((f:h)\leqslant (f:g)(g:h)\) 对于 \(f\in \mathcal{K}\) \(,g\in \mathcal{K}_{+}^{*}\) \(h\in \mathcal{K}_{+}^{*};\) 成立；
\end{enumerate}

\begin{enumerate}
\def\labelenumi{(\roman{enumi})}
\setcounter{enumi}{5}
\item
  \(0 < \frac{1}{(f_0:f)}\leqslant \frac{(f:g)}{(f_0:g)}\leqslant (f:f_0)\) 对于 \(f,f_0,g\) 属于 \(\mathcal{K}_+^*\) 成立；
\item
  设 \(f, f', h\) 属于 \(\mathcal{K}_+\)，且 \(h(s) \geqslant 1\) 在 \(f + f'\) 的支集上成立，再设 \(\varepsilon > 0\)；存在一个紧邻域 \(V\)，以 \(e\) 为中心，使得对每个 \(g \in \mathcal{K}_+^*(V)\)，
\end{enumerate}

\[
(f: g) + (f ^ {\prime}: g) \leqslant ((f + f ^ {\prime}): g) + \varepsilon (h: g).
\]

性质 (i)、(ii)、(iii) 显然。取 \(f \in \mathcal{K}\)、\(g \in \mathcal{K}_+^*\)；若 \(f \leqslant \sum_{i=1}^{n} c_i \gamma(s_i)g\)，且 \(c_i \geqslant 0\)，则 \(\sup f \leqslant \sum_{i=1}^{n} c_i g(s_i^{-1}s)\) 对某个 \(s \in \mathbf{G}\) 成立，因而 \(\sup f \leqslant \left( \sum_{i=1}^{n} c_i \right) \sup g\)，从而得到 (iv)。现在证明 (v)；取 \(f \in \mathcal{K}\)、\(g, h\) 属于 \(\mathcal{K}_+^*\)；若 \(f \leqslant \sum_{i=1}^{n} c_i \boldsymbol{\gamma}(s_i) g\) 且 \(g \leqslant \sum_{j=1}^{p} d_j \boldsymbol{\gamma}(t_j) h\)（\(c_i \geqslant 0\)、\(d_j \geqslant 0\)、\(s_i, t_j\) 属于 G），则 \(f \leqslant \sum_{i,j} c_i d_j \boldsymbol{\gamma}(s_i t_j) h\)，故 \((f: h) \leqslant \sum_{i,j} c_i d_j = \left( \sum_{i} c_i \right) \left( \sum_{j} d_j \right)\)；于是 \((f: h) \leqslant (f: g)(g: h)\)。把 (v) 一方面应用于 \(f_0, f, g\)，另一方面应用于 \(f, f_0, g\)，得到 (vi)。最后，设 \(f, f', h\) 属于 \(\mathcal{K}_+\)，且 \(h(s) \geqslant 1\) 在 \(f + f'\) 的支集上成立，再设 \(\varepsilon > 0\)。令 \(\mathbf{F} = f + f' + \frac{1}{2} \varepsilon h\)；函数 \(\varphi, \varphi'\) 分别与 \(f / F\) 和 \(f'/F\) 一致于 \(f + f'\) 的支集上，并在其外为零，属于 \(\mathcal{K}_+\)；对每个 \(\eta > 0\)，存在一个紧邻域 V，以 \(e\) 为中心，使得 \(|\varphi(s) - \varphi(t)| \leqslant \eta\) 且 \(|\varphi'(s) - \varphi'(t)| \leqslant \eta\) 在 \(s^{-1} t \in V\) 时成立。于是取 \(g \in \mathcal{K}_+^*(V)\)；对每个 \(s \in G\)，

\[
\varphi \cdot \boldsymbol {\gamma} (s) g \leqslant (\varphi (s) + \eta) \cdot \boldsymbol {\gamma} (s) g;
\]

这是因为，在 \(\pmb{\gamma}(s)g\) 为零的点，即在 \(s\mathbf{V}\) 之外，该不等式显然成立；而在 \(s\mathbf{V}\) 内，\(\varphi \leqslant \varphi(s) + \eta\)；同样，\(\varphi' \cdot \pmb{\gamma}(s)g \leqslant (\varphi'(s) + \eta) \cdot \pmb{\gamma}(s)g\)。现设数 \(c_1, \ldots, c_n\) 满足 \(\geqslant 0\)，并设 \(s_1, \ldots, s_n\) 为 \(\mathbf{G}\) 中的元素，使得 \(\mathbf{F} \leqslant \sum_{i=1}^{n} c_i \pmb{\gamma}(s_i)g\)；则

\[
f = \varphi \mathrm{F} \leqslant \sum_ {i = 1} ^ {n} c _ {i} \varphi \cdot \boldsymbol {\gamma} (s _ {i}) g \leqslant \sum_ {i = 1} ^ {n} c _ {i} (\varphi (s _ {i}) + \eta) \cdot \boldsymbol {\gamma} (s _ {i}) g
\]

对 \(f'\) 同理；因此

\[
(f: g) + (f ^ {\prime}: g) \leqslant \sum_ {i = 1} ^ {n} c _ {i} (\varphi (s _ {i}) + \varphi^ {\prime} (s _ {i}) + 2 \eta) \leqslant (1 + 2 \eta) \sum_ {i = 1} ^ {n} c _ {i}
\]

因为 \(\varphi +\varphi^{\prime}\leqslant 1\)。应用 F 的定义，再应用 (ii)、(iii) 和 (v)，可得

\[
\begin{array}{r l} & {(f: g) + (f ^ {\prime}: g) \leqslant (1 + 2 \eta) (\mathrm{F}: g) \leqslant} \\ & {\quad (1 + 2 \eta) \big [ \big ((f + f ^ {\prime}): g \big) + \frac {1}{2} \varepsilon (h: g) \big ] \leqslant} \\ & {\quad \big ((f + f ^ {\prime}): g \big) + \frac {1}{2} \varepsilon (h: g) + 2 \eta \big ((f + f ^ {\prime}): h \big) (h: g) + \varepsilon \eta (h: g)} \end{array}
\]

并且，若选择 \(\eta\) 使得 \(\eta\big[2\big((f + f'):h\big) + \varepsilon\big] \leqslant \frac{1}{2}\varepsilon\)，则得到 (vii)。

当 V 遍历 \(e\) 的紧邻域集合时，\(\mathcal{K}_{+}^{*}(V)\) 构成滤子 \(\mathfrak{B}\) 在 \(\mathcal{K}_{+}^{*}\) 上的一个基。令 \(\mathfrak{F}\) 为 \(\mathcal{K}_{+}^{*}\) 上比 \(\mathfrak{B}\) 更细的超滤子。另一方面，固定 \(f_{0} \in \mathcal{K}_{+}^{*}\)，并对 \(f \in \mathcal{K}_{+}^{*}\) 和 \(g \in \mathcal{K}_{+}^{*}\) 置

\[
\mathrm{I} _ {g} (f) = \frac {(f : g)}{(f _ {0} : g)}.
\]

由 (vi)，\(\lim_{g,\mathfrak{F}}\mathrm{I}_g(f) = \mathrm{I}(f)\) 存在于紧空间 \([1 / (f_0:f),(f:f_0)]\) 中。由 (iii)，\(\mathrm{I}(f + f')\leqslant \mathrm{I}(f) + \mathrm{I}(f')\)。由 (vii)，\(\mathrm{I}(f) + \mathrm{I}(f')\leqslant \mathrm{I}(f + f') + \varepsilon \mathrm{I}(h)\) 对所有 \(\varepsilon >0\) 成立，只要 \(h\) 满足 \(\geqslant 1\) 于 \(f + f'\) 的支集上；由此 \(\mathrm{I}(f + f') = \mathrm{I}(f) + \mathrm{I}(f')\)。根据 Ch. II, §2, No.~1, Prop. 2，I 可延拓为 \(\mathcal{K}\) 上的线性形式；该线性形式是 G 上的非零正测度，并由 (i) 左不变；这就是所求的左哈尔测度。转到对立群即可推出右哈尔测度的存在。

\begin{enumerate}
\def\labelenumi{\Alph{enumi})}
\setcounter{enumi}{1}
\tightlist
\item
  唯一性。--- 设 \(\mu\) 为左哈尔测度，\(\nu\) 为右哈尔测度。则 \(\check{\nu}\) 是左哈尔测度。我们将证明 \(\mu\) 与 \(\check{\nu}\) 成比例。这将证明任意两个左哈尔测度确实成比例。
\end{enumerate}

取 \(f \in \mathcal{K}\) 使得 \(\mu(f) \neq 0\)。由引理 1，按下式定义的函数 \(D_f\) 在 \(G\) 上

\[
\mathrm{D} _ {f} (s) = \mu (f) ^ {- 1} \int f (t ^ {- 1} s) d \nu (t)\tag{16}
\]

在 G 上连续。取 \(g \in \mathcal{K}\)。函数 \((s, t) \mapsto f(s)g(ts)\) 在 G × G 中连续且具有紧支集。由 Ch. III, §4, No.~1, Th. 2，

\[
\begin{array}{r l} & {\mu (f) \nu (g) = \Big (\int f (s) d \mu (s) \Big) \Big (\int g (t) d \nu (t) \Big)} \\ & {\qquad = \int d \mu (s) \int f (s) g (t s) d \nu (t) = \int d \nu (t) \int f (s) g (t s) d \mu (s)} \\ & {\qquad = \int d \nu (t) \int f (t ^ {- 1} s) g (s) d \mu (s)} \\ & {\qquad = \int g (s) \Big [ \int f (t ^ {- 1} s) d \nu (t) \Big ] d \mu (s) = \mu \big (g \cdot \mu (f) D _ {f} \big),} \end{array}\tag{17}
\]

从而

\[
\nu (g) = \mu (\mathrm{D} _ {f} \cdot g).\tag{18}
\]

首先，这证明 \(D_{f}\) 与 f 无关。对于 \(f' \in H\)，若满足 \(\mu(f') \neq 0\)，则有 \(D_{f} \cdot \mu = D_{f'} \cdot \mu\)，因此 \(D_{f} = D_{f'}\) 在 \(\mu\) 下局部几乎处处成立，进而处处成立，因为 \(\mathbf{D}_f\) 和 \(\mathbf{D}_{f'}\) 连续且 \(\mu\) 的支集为 \(\mathbf{G}\)。因此可记 \(\mathrm{D}_f = \mathrm{D}\)。公式 (16) 给出

\[
\mu (f) \mathrm{D} (e) = \stackrel {\vee} {\nu} (f).\tag{19}
\]

公式 (19) 可通过线性延拓到函数 \(f \in \mathcal{K}\)，其中满足 \(\mu(f) = 0\)。有 \(D(e) \neq 0\)，因为 \(\check{\nu} \neq 0\)。这就确立了 \(\mu\) 与 \(\check{\nu}\) 的比例关系。

推论。--- G 上每个左不变（分别地，右不变）测度都与某个左（分别地，右）哈尔测度成比例。

例。--- 1) 在加法群 R 上，勒贝格测度 dx 是哈尔测度（Ch. III, §1, No.~3, 例）。

\begin{enumerate}
\def\labelenumi{\arabic{enumi})}
\setcounter{enumi}{1}
\tightlist
\item
   对于每个函数 \(f \in \mathcal{K}(\mathbf{R}_+^*)\)，有（FRV, II, §1, 公式 (12)）
\end{enumerate}

\[
\int_ {0} ^ {+ \infty} \frac {f (x)}{x} d x = \int_ {0} ^ {+ \infty} \frac {f (t x)}{t x} t d x = \int_ {0} ^ {+ \infty} \frac {f (t x)}{x} d x
\]

对所有 t \textgreater{} 0，测度 \(x^{-1} dx\) 因而是乘法群 \(R_{+}^{*}\) 上的哈尔测度。

\begin{enumerate}
\def\labelenumi{\arabic{enumi})}
\setcounter{enumi}{2}
\tightlist
\item
  取 G 为圆环 \(\mathbf{T} = \mathbf{R} / \mathbf{Z}\)。令 \(\varphi\) 为从 \(\mathbf{R}\) 到 \(\mathbf{T}\) 的典范映射。对于 \(f \in \mathcal{K}(\mathbf{T})\)，函数 \(f \circ \varphi\) 在 \(\mathbf{R}\) 上连续且以 1 为周期，积分
\end{enumerate}

\[
\mathrm{I} (f) = \int_ {a} ^ {a + 1} f (\varphi (x)) d x
\]

与 \(a \in \mathbf{R}\) 的选择无关；它显然关于平移不变；因此它定义了 \(\mathbf{T}\) 上的哈尔测度。通过结构传递，由此得到 \(I(f) = \int_{a}^{a + 1} f(e^{2\pi it}) dt\) 是复数绝对值为 1 的乘法群 \(\mathbf{U}\) 上的哈尔测度（GT, VIII, §2, No.~1）。

命题 2. --- 设 G 是局部紧群，\(\mu\) 是 G 上的左或右哈尔测度。G 为离散的充要条件是 \(\mu(\{e\}) > 0\)。G 为紧的充要条件是 \(\mu^{*}(G) < +\infty\)。

这些条件显然是必要的。证明它们的充分性。令 V 为 e 的紧邻域。若 \(\mu(\{e\}) > 0\)，则由于 \(\mu(V) < +\infty\)，V 是有限集；由于 G 是豪斯多夫空间，因而 G 是离散的。假设 \(\mu^{*}(G)<+\infty\)，并且例如 \(\mu\) 左不变。考虑集合 \(\mathscr{E}\)，它由 G 的有限子集 \(\{s_{1},\ldots,s_{n}\}\) 组成，且满足 \(s_{i}V\cap s_{j}V=\varnothing\) 当 \(i\neq j\) 时；则

\[
n \mu (\mathrm{V}) = \mu (s _ {1} \mathrm{V} \cup \dots \cup s _ {n} \mathrm{V}) \leqslant \mu^ {*} (\mathrm{G}),
\]

因此 \(n \leqslant \mu^{*}(\mathrm{G})/\mu(\mathrm{V})\)。于是可以在 E 中选取一个极大元 \(\{s_{1},\ldots,s_{n}\}\)。然后，对于每个 \(s \in G\)，存在某个 i 使得 \(sV \cap s_{i}V \neq \varnothing\)，从而有 \(s \in s_{i}VV^{-1}\)。因此 G 是紧集 \(s_{i}VV^{-1}\) 的并，故为紧的。

